\documentclass[10pt,a4paper]{amsart}
\usepackage[margin=19mm]{geometry}
\usepackage{amsmath,amssymb,mathtools}
\usepackage[scr=rsfs,cal=euler]{mathalfa}
\usepackage{xfrac}
\usepackage[unicode,hidelinks,bookmarks=false]{hyperref}
\newcommand{\ie}{i.e.\ }
\newcommand{\cf}{cf.\ }
\newcommand{\resp}{respectively\ }
\newcommand{\Subsection}{\subsection}
\providecommand{\nobreakdash}{\nobreak\hbox{-}}
\setlength{\emergencystretch}{2em}
\multlinegap0pt
\newtheorem{comment}{Comment}
\usepackage{amssymb}
\usepackage{enumerate}
\usepackage{tikz}
\usepackage{xcolor}
\newtheorem{lemma}{Lemma}
\newtheorem{cor}{Corollary}
\newcommand{\Ma}{\color{blue}}
\newcommand\N{\mathbb N}
\renewcommand\O\Omega
\newcommand{\Ying}{\color{red}}
\newcommand\e\varepsilon
\title{Qualitative properties of eigenfunctions in domains with small holes}
\author{Laura Abatangelo, Massimo Grossi, Ying Li}
\date{Source: arXiv:2607.29487v1 (CC BY 4.0)}
\begin{document}
\maketitle

\noindent\emph{Source and licence.} Qualitative properties of eigenfunctions in domains with small holes. Version arXiv:2607.29487v1 (CC BY 4.0), distributed by arXiv under the Creative Commons Attribution 4.0 International licence, http://creativecommons.org/licenses/by/4.0/, as recorded in the arXiv licence metadata for this exact version and shown on the abstract page https://arxiv.org/abs/2607.29487v1. Full licence notice: \texttt{licenses/arxiv-2607.29487-CC-BY-4.0.txt}.\par\smallskip

\noindent\textit{Editorial scope.} Counted unit: the article's lemma l:improvedL2Ve (source0.tex lines 1137--1176). The article's complete original proof of it is delivered with it (source0.tex lines 1177--1415, one \texttt{proof} environment of 239 lines). No mathematics is written, repaired or re-derived: the statement and the proof are the article's own lines, in the article's own notation, and no retained line is substituted or re-flowed.  Retained uncounted as the source-local context the proof needs: lines 177--182; lines 3531--3537; lines 3647--3655.  Nothing was omitted from any retained span, so the package carries no omission record.  No retained line contains a citation, so the wrapper carries no bibliography and no bibliography entry.  No retained line contains a figure environment, an image-inclusion command or drawing code, so no image is delivered and the package declares no asset dependency.  Editorial formatting only, in addition: an AMS \texttt{amsart} preamble that restates the article's own declarations and macros used by the retained lines, namely the environments \texttt{lemma}, \texttt{cor} on separate counters, together with the standard packages those lines need.  The retained environments print in this document as Lemma 1, Corollary 1 in order of appearance, whereas the article prints the same items with the numbering of its own full document.  The complete native source of the article as distributed in the arXiv e-print for this exact version is bundled unchanged as \texttt{sources/206491/source0.tex}.
	\begin{equation}\label{eq:eigenvalueproblem}
		\begin{cases}
			-\Delta u=\lambda u~&\mbox{in}\ \Omega,\\
			u=0~&\mbox{on}\ \partial\Omega.
		\end{cases}
	\end{equation}

       \begin{lemma}\label{lem10102025-1}
             Let $N\ge 2$.   We claim that, for any  $k\ge2$,
					\begin{equation}\label{eq1609-1}
						\sum_{i_1,i_2,\cdots,i_k=1}^{N}\frac{\partial^k u(x)}{\partial x_{i_1} \partial x_{i_2}\cdots \partial x_{i_k}}Q_{i_1,i_2,\cdots,i_k}(x,y)=	\sum_{i_1,i_2,\cdots,i_{k-2}=1}^{N}\frac{\partial^{k-2}(\Delta u) (x)}{\partial x_{i_1} \partial x_{i_2}\cdots \partial x_{i_{k-2}}} \widetilde{Q}_{i_1,i_2,\cdots,i_{k-2}}(x,y),
					\end{equation}
					where $\widetilde{Q}_{i_1,i_2,\cdots,i_{k-2}}(x,y)$ is a polynomial of degree k depending on $i_1,i_2,\cdots,i_{k-2}$.
            \end{lemma}

                 \begin{cor}\label{Cor-Q}
                 Assume that $k\in\N$ is the vanishing order of $u$ at $P$.
                     If $u$ is an eigenfunction of Dirichlet Laplacian on $\O$ satisfying \eqref{eq:eigenvalueproblem}, then
                      \begin{equation}\label{eq2904-2}
					\sum_{|\alpha|=k}\frac1{\alpha!}
					\frac{\partial^ku(P)}{\partial x_1^{\alpha_1}\cdots\partial x_N^{\alpha_N}}Q_\alpha(x,P)=0\ \ \hbox{and}\ \ 	\sum_{|\beta|=k+1}\frac1{\beta!}
					\frac{\partial^{k+1}u(P)}{\partial x_1^{\beta_1}\cdots\partial x_N^{\beta_N}}Q_\beta(x,P)=0.
				\end{equation}   
                 \end{cor}  

\begin{lemma}\label{l:improvedL2Ve}
 For $N\ge 3$, let 	$A(N,0)=N(N-2)\omega_N$ and 
			$A(N,k)=\frac{N(N-2)\omega_N}{(N-2)N\cdots(N-4+2k)}$, for any $k\ge1$. 
			We have that
				\begin{equation}\label{V1}
				\!\!\!\!\!\!\!\!\!\!\!\!V_{\e,u}(x)=A(N,k)\e^{N-2+2k}
					\sum_{|\alpha|=k}\frac1{\alpha!}
					\left(\frac{\partial^ku(P)}{\partial x_1^{\alpha_1}\cdots\partial x_N^{\alpha_N}}+O(\e)\right)
					\left(\frac{\partial^kG_{\Omega}(x,P)}{\partial y_1^{\alpha_1}\cdots\partial y_N^{\alpha_N}}+O(\e)\right)+O(\e^{N+k})G_{\Omega}(x,P),
				\end{equation}
              %  {\Ma One question: in the final part of the proof we pointed out that if $N=3$ and $k=0$ the lower order term is $O(\e^2)$ instead of $O(\e^3)$. Maybe we have to insert it } {\Ying Yes, the lower order term is $O(\e^2).$ But in the first term, it contains  $O(\e^2).$ For example, $\e^{N-2+2k}\left(\frac{\partial^ku(P)}{\partial x_1^{\alpha_1}\cdots\partial x_N^{\alpha_N}}+O(\e)\right)\left(\frac{\partial^kG_{\Omega}(x,P)}{\partial y_1^{\alpha_1}\cdots\partial y_N^{\alpha_N}}+O(\e)\right)=\e^{N-2+2k}\frac{\partial^ku(P)}{\partial x_1^{\alpha_1}\cdots\partial x_N^{\alpha_N}}\frac{\partial^kG_{\Omega}(x,P)}{\partial y_1^{\alpha_1}\cdots\partial y_N^{\alpha_N}}+O(\e^{N-1+2k})\frac{\partial^kG_{\Omega}(x,P)}{\partial y_1^{\alpha_1}\cdots\partial y_N^{\alpha_N}}+O(\e^{N-1+2k})$ } {\Ma Ok, I did not see the term $O(\e^2)$}
			as $\e\to0$,	uniformly for $x\in\O_\e$.
			Furthermore, we obtain
			\begin{equation}\label{V3}
				\begin{aligned}
				&\hbox{For $k=0$,}\quad
					||V_{\e,u}||_{L^p(\O_\e)}=
					\begin{cases}
					\displaystyle	O(\e^{N-2}),  &if\ 1\le p<\frac N{N-2},\\
						O(\e^{N-2}|\log\e|), &if\  p=\frac{N}{N-2},\\
						O(\e^{\frac{N}{p}}),   &if\  p>\frac{N}{N-2}.\\
					\end{cases}\\
				&\hbox{For $k=1$,}\quad
					||V_{\e,u}||_{L^p(\O_\e)}=
					\begin{cases}
						O(\e^{N}),  &if\ 1\le p<\frac N{N-1},\\
						O(\e^{N}|\log\e|), &if\  p=\frac{N}{N-1},\\
						O(\e^{\frac{N}{p}+1}),   &if\  p>\frac{N}{N-1}.\\
					\end{cases}\\
			&	\hbox{For $k= 2$,}\quad
					||V_{\e,u}||_{L^p(\O_\e)}=
					\begin{cases}
						O(\e^{N+2}|\log\e|), &if\  p=1,\\
						O(\e^{\frac{N}{p}+2}),   &if\  p>1.\\
					\end{cases}\\
		&	\hbox{For $k\ge 3$,}\quad
					||V_{\e,u}||_{L^p(\O_\e)}=	O(\e^{\frac{N}{p}+k}),\quad if \ p\ge 1.
				\end{aligned}
			\end{equation}
		\end{lemma}

		\begin{proof}[\bf Proof of \eqref{V1}.]
			Let us show that the function 
            \begin{equation}\label{eqwe}
				\begin{aligned}
					w_\e=&A(N,k)\e^{N-2+2k}	\sum_{|\alpha|=k}\frac1{\alpha!}
					\frac{\partial^ku(P)}{\partial x_1^{\alpha_1}\cdots\partial x_N^{\alpha_N}}
					\frac{\partial^kG_{\Omega}(x,P)}{\partial y_1^{\alpha_1}\cdots\partial y_N^{\alpha_N}}\\
					&+A(N,k+1)\e^{N+2k}	\sum_{|\alpha|=k+1}\frac1{\alpha!}
					\frac{\partial^{k+1}u(P)}{\partial x_1^{\alpha_1}\cdots\partial x_N^{\alpha_N}}
					\frac{\partial^{k+1}G_{\Omega}(x,P)}{\partial y_1^{\alpha_1}\cdots\partial y_N^{\alpha_N}}
				\end{aligned}      
            \end{equation}
			verifies
			\begin{equation}
				\begin{cases}
					\Delta w_\e=0&in\ \O_\e,\\
					w_\e=0&on\ \partial\O.
				\end{cases}
			\end{equation}
			  Let $Q_\alpha(x,P)$ be a polynomial of degree $k$ depending on $\alpha=(\alpha_1,\cdots,\alpha_N)$ such that \eqref{eq0712-1} holds,
				\begin{equation}\label{eq0712-1}
					\frac{\partial^k\left(|x-y|^{2-N}\right)}{\partial y_1^{\alpha_1}\cdots\partial y_N^{\alpha_N}}\Bigg|_{y=P}=C(N,k)\frac{(x_1-P_1)^{\alpha_1}\cdots(x_N-P_N)^{\alpha_N}+Q_\alpha(x,P)}
					{|x-P|^{N-2+2k}},
				\end{equation}
			where  $C(N,0)=1$, $C(N,k)=(N-2)N\cdots(N-4+2k)$ for any $k\ge 1$. 
            \begin{comment}
                {\Ma Are you sure about the value of $C(N,k)$? It seems that for $k=1$ we should have $C(N,1)=N(N-2)$. {\Ying (for $k=1$, $C(N,1)=(N-2), C(N,2)=(N-2)N$, the value of $C(N,k)$ looks right, but this expression is probably misleading)}} {\Ma I agree that for for $k=1$, $C(N,1)=(N-2)and C(N,2)=(N-2)N$, but if you write as above it is not correct}\\
            {\Ying you mean that \begin{align*}
			   &\frac{\partial}{\partial y_1}\frac{(x_1-y_1)^{\alpha_1}\cdots(x_N-y_N)^{\alpha_N}}{|x-y|^{N-2+2k}}\\
               =&\frac{(N-2+2k)(x_1-y_1)^{\alpha_1+1}(x_2-y_2)^{\alpha_2}\cdots(x_N-y_N)^{\alpha_N}-\alpha_1(x_1-y_1)^{\alpha_1-1}(x_2-y_2)^{\alpha_2}\cdots(x_N-y_N)^{\alpha_N}|x-y|^2}{|x-y|^{N+2k}}\\
               =&\frac{(N-2+2k-\alpha_1)(x_1-y_1)^{\alpha_1+1}(x_2-y_2)^{\alpha_2}\cdots(x_N-y_N)^{\alpha_N}-\alpha_1(x_1-y_1)^{\alpha_1-1}(x_2-y_2)^{\alpha_2}\cdots(x_N-y_N)^{\alpha_N}((x_2-y_2)^2+\cdots+(x_N-y_N)^2)}{|x-y|^{N+2k}} 
			\end{align*}
            I just throw the term $\alpha_1(x_1-y_1)^{\alpha_1-1}(x_2-y_2)^{\alpha_2}\cdots(x_N-y_N)^{\alpha_N}|x-y|^2$ into $Q_\alpha(x,y)$, i.e., I just keep the term $(N-2+2k)(x_1-y_1)^{\alpha_1+1}(x_2-y_2)^{\alpha_2}\cdots(x_N-y_N)^{\alpha_N}$, the other polynomials are thrown into $Q_\alpha(x,y)$. I will check it and Lemma \ref{lem10102025-1} again.\\
            } and $Q_\alpha(x,P)$ is a polynomial of degree $k$ depending on $\alpha_1,\cdots,\alpha_N$. Observe that 
                \begin{equation}
                   \sum_{|\alpha|=k}\frac1{\alpha!}
					\frac{\partial^ku(P)}{\partial x_1^{\alpha_1}\cdots\partial x_N^{\alpha_N}}Q_\alpha(x,P) =\frac{1}{k!}\sum_{i_1,i_2,\cdots,i_k=1}^{N}\frac{\partial^k u(P)}{\partial x_{i_1} \partial x_{i_2}\cdots \partial x_{i_k}}Q_{i_1,i_2,\cdots,i_k}(x,P),
                \end{equation}
                where $Q_{i_1,i_2,\cdots,i_k}(x,y)$ {\Ma Ying, here you have to say some words about this notation. For the reader I do not know how clear is. Maybe it is better to keep the notation with $\alpha$?} {\Ying ($Q_{i_1,\cdots,i_k}$ is remaining term of taking derivative of $|x-y|^{N-2}$ with respect to $y_{i_1},\cdots,y_{i_k}$ I am thinking a good definition of $Q_{i_1,\cdots,i_k}$) } is a polynomial of degree $k$ depending on $i_1,i_2,\cdots,i_k$, defined by
				\begin{equation}\label{eq10102025-4}
					\frac{\partial^k\left(|x-y|^{2-N}\right)}{\partial y_{i_1}\cdots\partial y_{i_k}}=C(N,k)\frac{(x_{i_1}-y_{i_1})\cdots(x_{i_k}-y_{i_k})+Q_{i_1,i_2,\cdots,i_k}(x,y)}
					{|x-P|^{N-2+2k}}.
				\end{equation}
                By direct computations, we have $Q_{i}(x,y)=0$ {\Ma Sorry, but this notation here is not clear. You mean $Q_{i}(x,y)=0$ for $k=1$? In this case you mean that only one of the indices $i_1,i_2,\cdots,i_k$ is not zero? {\Ying ($Q_{i_1}(x,y)=0$, for any $i_1=1,\cdots,N$. It is better to use the uniform indices which I will revise.)}} {\Ma OK, I think that the definition of Q is quite clear, there is some troubles to understand the notation $i_1,i_2,\cdots,i_k$} {\Ying For example, $N=2$,  $k=3$, then $i_1,i_2,i_3=1,2$,
                $$
                \frac{\partial^3 u}{\partial y_{i_1}\partial y_{i_2}\partial y_{i_3}}=\begin{cases}
                    \frac{\partial^3 u}{\partial y_1^3},&if \ (i_1,i_2,i_3)=(1,1,1)\\
                     \frac{\partial^3 u}{\partial y_1^2\partial y_2},&if \ (i_1,i_2,i_3)=(1,1,2)\\
                      \frac{\partial^3 u}{\partial y_1^2\partial y_2},&if \ (i_1,i_2,i_3)=(1,2,1)\\
                       \frac{\partial^3 u}{\partial y_1\partial y_2^2},&if \ (i_1,i_2,i_3)=(1,2,2)\\
                        \frac{\partial^3 u}{\partial y_1^2\partial y_2},&if \ (i_1,i_2,i_3)=(2,1,1)\\
                        \frac{\partial^3 u}{\partial y_1\partial y_2^2},&if \ (i_1,i_2,i_3)=(2,1,2)\\
                        \frac{\partial^3 u}{\partial y_1\partial y_2^2},&if \ (i_1,i_2,i_3)=(2,2,1)\\
                         \frac{\partial^3 u}{\partial y_2^3},&if \ (i_1,i_2,i_3)=(2,2,2)\\
                \end{cases}
                $$} and $Q_{i,j}(x,y)=-\frac{1}{N}\delta_i^j|x-y|^2$ {\Ma Same comment, we are in the case $k=2$?} {\Ying ($Q_{i_i,i_2}(x,y)=-\frac{1}{N}\delta_{i_1}^{i_2}|x-y|^2$, $i_1,i_2=1,\cdots,N$)}
                Then for $k=0$ {\Ma $k=1$?}{\Ying (This sentence is to say for $k=0$ and $k=1$, \eqref{eq1704-1} and \eqref{eq1704-2} hold) Because, in Lemma \ref{lem10102025-1} holds for $k\ge 2$. Maybe I can write a Corollary about \eqref{eq1704-1} and \eqref{eq1704-2} separately) }, we see that $\sum_{i=1}^N\frac{\partial u(P)}{\partial x_i}Q_i(x,P)=0$. For $k=1$ {\Ma $k=2$?}, we  have
                $
\sum_{i=1}^N\frac{\partial u(P)}{\partial x_i}Q_i(x,P)=0
                $
                and
                \[
                \sum_{i,j=1}^N
				\frac{\partial^2 u(P)}{\partial x_i \partial x_j} Q_{i,j}(x,P)=
				-\frac{1}{N}\sum_{i,j=1}^N
				\frac{\partial^2 u(P)}{\partial x_i \partial x_j} \delta_i^j |x - P|^2 = -\frac{1}{N}|x - P|^2 \Delta u(P)=\frac{\lambda}{N}|x-P|^2u(P)=0 .
				\]
                 \end{comment}
            Since the vanishing order of $u$ at $P$ is $k$, it follows from Corollary \ref{Cor-Q} in the Appendix that for any $k\ge 0$
             \begin{equation}\label{eq1704-1}
					\sum_{|\alpha|=k}\frac1{\alpha!}
					\frac{\partial^ku(P)}{\partial x_1^{\alpha_1}\cdots\partial x_N^{\alpha_N}}Q_\alpha(x,P)=0,
				\end{equation}   
                and 
                \begin{equation}\label{eq1704-2}
					\sum_{|\alpha|=k+1}\frac1{\alpha!}
					\frac{\partial^{k+1}u(P)}{\partial x_1^{\alpha_1}\cdots\partial x_N^{\alpha_N}}Q_\alpha(x,P)=0.
				\end{equation}   
            \begin{comment}
				\begin{align*}
					\!\!\!\!\!\!\!\!\!\!\!\!\!\!\!\!\!\!\!\!\!
                    w_\e(x)=&\frac{1}{(N-2)N\cdots(N-4+2k)}\e^{N-2+2k}
					\sum_{|\alpha|=k}\frac1{\alpha!}
					\frac{\partial^ku(P)}{\partial x_1^{\alpha_1}\cdots\partial x_N^{\alpha_N}}
					\frac{\partial^k\left(|x-y|^{2-N}+CH_{\Omega}(x,y)\right)}{\partial y_1^{\alpha_1}\cdots\partial y_N^{\alpha_N}}\Bigg|_{y=P}\\
                    &+\frac{1}{(N-2)N\cdots(N-2+2k)}\e^{N+2k}
					\sum_{|\alpha|=k+1}\frac1{\alpha!}
					\frac{\partial^{k+1}u(P)}{\partial x_1^{\alpha_1}\cdots\partial x_N^{\alpha_N}}
					\frac{\partial^{k+1}\left(|x-y|^{2-N}+CH_{\Omega}(x,y)\right)}{\partial y_1^{\alpha_1}\cdots\partial y_N^{\alpha_N}}\Bigg|_{y=P}\\
                    =&\frac{\e^{N-2+2k}}{|x-P|^{N-2+2k}}\sum_{|\alpha|=k}\frac1{\alpha!}
					\frac{\partial^ku(P)}{\partial x_1^{\alpha_1}\cdots\partial x_N^{\alpha_N}}(x_1-P_1)^{\alpha_1}\cdots(x_N-P_N)^{\alpha_N}+O(\e^{N-2+2k})\\
                    &+\frac{\e^{N+2k}}{|x-P|^{N+2k}}\sum_{|\alpha|=k+1}\frac1{\alpha!}
					\frac{\partial^{k+1}u(P)}{\partial x_1^{\alpha_1}\cdots\partial x_N^{\alpha_N}}(x_1-P_1)^{\alpha_1}\cdots(x_N-P_N)^{\alpha_N}+O(\e^{N+2k}).\\
				\end{align*}
                     \end{comment}
                     Therefore, for any $k\ge 0$,
                     by \eqref{eqwe} we get
                \begin{align*}
					\!\!\!\!\!\!\!\!\!\!\!\!\!\!\!\!\!\!\!\!\!
                    w_\e(x)=&A(N,k)\e^{N-2+2k}
					\sum_{|\alpha|=k}\frac1{\alpha!}
					\frac{\partial^ku(P)}{\partial x_1^{\alpha_1}\cdots\partial x_N^{\alpha_N}}
					\frac{\partial^k\left(\frac1{N(N-2)\omega_N}|x-y|^{2-N}+H_{\Omega}(x,y)\right)}{\partial y_1^{\alpha_1}\cdots\partial y_N^{\alpha_N}}\Bigg|_{y=P}\\
                    &+A(N,k+1)\e^{N+2k}
					\sum_{|\alpha|=k+1}\frac1{\alpha!}
					\frac{\partial^{k+1}u(P)}{\partial x_1^{\alpha_1}\cdots\partial x_N^{\alpha_N}}
					\frac{\partial^{k+1}\left(\frac1{N(N-2)\omega_N}|x-y|^{2-N}+H_{\Omega}(x,y)\right)}{\partial y_1^{\alpha_1}\cdots\partial y_N^{\alpha_N}}\Bigg|_{y=P}\\
                    =&\frac{\e^{N-2+2k}}{|x-P|^{N-2+2k}}\sum_{|\alpha|=k}\frac1{\alpha!}
					\frac{\partial^ku(P)}{\partial x_1^{\alpha_1}\cdots\partial x_N^{\alpha_N}}(x_1-P_1)^{\alpha_1}\cdots(x_N-P_N)^{\alpha_N}+O(\e^{N-2+2k})\\
                    &+\frac{\e^{N+2k}}{|x-P|^{N+2k}}\sum_{|\alpha|=k+1}\frac1{\alpha!}
					\frac{\partial^{k+1}u(P)}{\partial x_1^{\alpha_1}\cdots\partial x_N^{\alpha_N}}(x_1-P_1)^{\alpha_1}\cdots(x_N-P_N)^{\alpha_N}+O(\e^{N+2k}).\\
				\end{align*}
                For $|x-P|=\e$, we have
                \begin{equation}\label{eq11102025-1}
                    \begin{aligned}
                    w_\e(x)=&\sum_{|\alpha|=k}\frac1{\alpha!}
					\frac{\partial^ku(P)}{\partial x_1^{\alpha_1}\cdots\partial x_N^{\alpha_N}}(x_1-P_1)^{\alpha_1}\cdots(x_N-P_N)^{\alpha_N}\\
                    &+\sum_{|\alpha|=k+1}\frac1{\alpha!}
					\frac{\partial^{k+1}u(P)}{\partial x_1^{\alpha_1}\cdots\partial x_N^{\alpha_N}}(x_1-P_1)^{\alpha_1}\cdots(x_N-P_N)^{\alpha_N}+O(\e^{N-2+2k}).
                    \end{aligned}
                \end{equation}
%And it is easy to see that  \eqref{eq11102025-1} holds for $k\ge0$.
Since $V_{\e,u}(x)=u(x)$ on $\partial B(P,\e)$, by using Taylor expansion for $x\in\partial B(P,\e)$ and combining with \eqref{eq11102025-1}, we obtain
\begin{equation}\label{eq1901-1}
			\begin{aligned}
				&V_{\e,u}(x)-w_\e(x)=u(x)-w_\e(x)\\
				=&\sum_{|\alpha|=k}\frac1{\alpha!}
					\frac{\partial^ku(P)}{\partial x_1^{\alpha_1}\cdots\partial x_N^{\alpha_N}}(x_1-P_1)^{\alpha_1}\cdots(x_N-P_N)^{\alpha_N}\\
                    &+\sum_{|\alpha|=k+1}\frac1{\alpha!}
					\frac{\partial^{k+1}u(P)}{\partial x_1^{\alpha_1}\cdots\partial x_N^{\alpha_N}}(x_1-P_1)^{\alpha_1}\cdots(x_N-P_N)^{\alpha_N}+O(\e^{k+2})\\
				&-\sum_{|\alpha|=k}\frac1{\alpha!}
					\frac{\partial^ku(P)}{\partial x_1^{\alpha_1}\cdots\partial x_N^{\alpha_N}}(x_1-P_1)^{\alpha_1}\cdots(x_N-P_N)^{\alpha_N}\\
                    &-\sum_{|\alpha|=k+1}\frac1{\alpha!}
					\frac{\partial^{k+1}u(P)}{\partial x_1^{\alpha_1}\cdots\partial x_N^{\alpha_N}}(x_1-P_1)^{\alpha_1}\cdots(x_N-P_N)^{\alpha_N}+O(\e^{N-2+2k})\\
                    =&O(\e^{k+2})+O(\e^{N-2+2k}).
			\end{aligned}
            \end{equation}
		If we directly estimate $V_{\e,u}-w_\e$ by the maximum principle for the harmonic function, the estimation is a little rough. So we use the Green function.\\ On $|x-P|=\e$, $(\e^{N+k}+\e^{2(N-2)+2k})G_{\Omega}(x,P)=\frac1{N(N-2)\omega_N}(\e^{k+2}+\e^{N-2+2k})+O(\e^{N+k})+O(\e^{2(N-2)+2k})$. Hence we can select a large $C>0$ such that   
				\begin{equation}
					\begin{cases}
						\Delta\Big[V_{\e,u}-w_\e-C(\e^{N+k}+\e^{2(N-2)+2k})G_{\Omega}(x,P)\Big]=0&in\ \O_\e,\\
						V_{\e,u}-w_\e-C(\e^{N+k}+\e^{2(N-2)+2k})G_{\Omega}(x,P)=0&on\ \partial\O,\\
						V_{\e,u}-w_\e-C(\e^{N+k}+\e^{2(N-2)+2k})G_{\Omega}(x,P)\le 0&on\ |x-P|=\e,
					\end{cases}
				\end{equation}
				and so by the maximum principle for harmonic function we get
				\begin{equation}
					V_{\e,u}-w_\e \le C(\e^{N+k}+\e^{2(N-2)+2k})G_{\Omega}(x,P)\quad\hbox{in }\O_\e,
				\end{equation}
				and since the same holds for the reverse inequality, we deduce that \begin{equation}\label{eq-0922-1}
					|	V_{\e,u}-w_\e |\leq C(\e^{N+k}+\e^{2(N-2)+2k})G_{\Omega}(x,P)\quad\hbox{in }\O_\e.
				\end{equation}
                Actually,
                 \begin{align*}
					&A(N,k+1) \e^{N+2k}	\sum_{|\alpha|=k+1}\frac1{\alpha!}
					\frac{\partial^{k+1}u(P)}{\partial x_1^{\alpha_1}\cdots\partial x_N^{\alpha_N}}
					\frac{\partial^{k+1}G_{\Omega}(x,P)}{\partial y_1^{\alpha_1}\cdots\partial y_N^{\alpha_N}}\\
                    =&\frac{\e^{N+2k}}{|x-P|^{N+2k}}\sum_{|\alpha|=k+1}\frac1{\alpha!}
					\frac{\partial^{k+1}u(P)}{\partial x_1^{\alpha_1}\cdots\partial x_N^{\alpha_N}}(x_1-P_1)^{\alpha_1}\cdots(x_N-P_N)^{\alpha_N}+O(\e^{N+2k})\\
                    =& O(\e)\frac{\e^{N-2+2k}}{|x-P|^{N-2+2k}}\sum_{|\alpha|=k}(x_1-P_1)^{\alpha_1}\cdots(x_N-P_N)^{\alpha_N}+O(\e^{N+2k})\\
                    =&O(\e)\e^{N-2+2k}	\sum_{|\alpha|=k}
					\frac{\partial^k\Big(G_{\Omega}(x,P)-CH_{\Omega}(x,P)\Big)}{\partial y_1^{\alpha_1}\cdots\partial y_N^{\alpha_N}}+O(\e^{N+2k})\\
                    =&O(\e)\e^{N-2+2k}	\sum_{|\alpha|=k}
					\frac{\partial^kG_{\Omega}(x,P)}{\partial y_1^{\alpha_1}\cdots\partial y_N^{\alpha_N}}+O(\e^{N-1+2k}).\\
                \end{align*}
                This implies that 
                 \begin{equation}\label{eq11102025-2}
                 \begin{aligned}
                   V_{\e,u}(x)=&A(N,k)\e^{N-2+2k}
					\sum_{|\alpha|=k}\frac1{\alpha!}
					\left(\frac{\partial^ku(P)}{\partial x_1^{\alpha_1}\cdots\partial x_N^{\alpha_N}}
                    +O(\e)\right)
					\left(\frac{\partial^kG_{\Omega}(x,P)}{\partial y_1^{\alpha_1}\cdots\partial y_N^{\alpha_N}}
                    +O(\e)\right)\\
                    &+O(\e^{N+k}+\e^{2(N-2)+2k})G_{\Omega}(x,P).  
                 \end{aligned}
				\end{equation}
                As final observation, we can see that $2(N-2)+2k<N+k$ if and only if $N=3$ and $k=0$. 
               In this case, the term $O(\e^2)G_\O(x,P)$ can be absorbed in the main term $A(3,0)\e (u(P)+O(\e))(G_\O(x,P)+O(\e))$.
                Except in this case, the term $O(\e^{2(N-2)+2k})$
                 is of lower order than $O(\e^{N+k})$, and therefore we have
   \begin{equation}\label{eq11102025-3}
					V_{\e,u}(x)=A(N,k)\e^{N-2+2k}
					\sum_{|\alpha|=k}\frac1{\alpha!}
					\left(\frac{\partial^ku(P)}{\partial x_1^{\alpha_1}\cdots\partial x_N^{\alpha_N}}+O(\e)\right)
					\left(\frac{\partial^kG_{\Omega}(x,P)}{\partial y_1^{\alpha_1}\cdots\partial y_N^{\alpha_N}}+O(\e)\right)+O(\e^{N+k})G_{\Omega}(x,P).
				\end{equation}
				{\bf Proof of \eqref{V3} .}\\
                By \eqref{eq11102025-3}, we have
				\begin{equation}
                \begin{aligned}
					||V_{\e,u}||_{L^p(\O_\e)}^p\le& C\e^{(N-2+2k)p}\int_{\O_\e}\sum_{|\alpha|=k}\Big|\frac{\partial^kG_{\Omega}(x,P)}{\partial y_1^{\alpha_1}\cdots\partial y_N^{\alpha_N}}\Big|^p +C\e^{(N+k)p}
					\int_{\O_\e}|G_{\Omega}(x,P)|^p\\
                    \le & C\e^{(N-2+2k)p}\int_{\O_\e} \frac{1}{|x-P|^{(N-2+k)p}}+C\e^{(N+k)p}
					\int_{\O_\e}\frac1{|x-P|^{(N-2)p}}=:I+J.\\
                     \end{aligned}
				\end{equation}
				 We estimate the terms $I$ and $J$ separately, the estimate \eqref{V3} it will follow from the combination of the two. For $k=0$,
				\begin{equation}
					I=
					\begin{cases}
						O(\e^{(N-2)p}),  &if\ 1\le p<\frac N{N-2},\\
						O(\e^{(N-2)p}|\log\e|), &if\  p=\frac{N}{N-2},\\
						O(\e^{N}),   &if\  p>\frac{N}{N-2}.\\
					\end{cases}
				\end{equation}
				For $k=1$,
				\begin{equation}
					I=
					\begin{cases}
						O(\e^{Np}),  &if\ 1\le p<\frac N{N-1},\\
						O(\e^{Np}|\log\e|), &if\  p=\frac{N}{N-1},\\
						O(\e^{N+p}),   &if\  p>\frac{N}{N-1}.\\
					\end{cases}
				\end{equation}
				For $k= 2$,
				\begin{equation}
					I=
					\begin{cases}
						O(\e^{(N+2)p}|\log\e|), &if\  p=1,\\
						O(\e^{N+2p}),   &if\  p>1.\\
					\end{cases}
				\end{equation}
				For $k\ge 3$,
				\begin{equation}
					I=	O(\e^{N+kp}),\quad if \ p\ge 1.
				\end{equation}
                For $k\ge 0$, 
				\begin{equation}
					J=
					\begin{cases}
						O(\e^{(N+k)p}),&if\ 1\le p<\frac N{N-2},\\
						O(\e^{(N+k)p}|\log\e|),&if\  p=\frac N{N-2},\\
						O(\e^{N+(k+2)p}),&if\ p>\frac N{N-2}.
					\end{cases}
				\end{equation}
               % By comparing $I$ and $J$ for each $k\ge 0$, we can see that $J\le CI$. 
                Then combining the above equations, we complete the proof of  \eqref{V3}.
				\end{proof} 	   	 

\end{document}
